% 単体ビルド: VS Code でこのファイルを開いてビルド，または
%   latexmk -cd ja/lessons/01-using-this-template.tex
\documentclass[../main.tex]{subfiles}
\begin{document}

\chapter{このテンプレートの使い方}
\lessoninfo{
  video={Lesson 1 (Introduction)},
  textbook={H\&P \cite{hennessy2017} §1.1--1.3},
  keywords={サイドノート, 索引, 二言語ビルド},
}

この章はクラスが提供する全マクロの見本である．実際の講義ノートを書き
始めるときは削除してよいが，\texttt{02-lesson-skeleton.tex} は新しい章の
雛形として残しておくとよい．

\section{余白ノート}

右側の余白は，本文の流れを妨げずに補足したい内容のための領域である．
番号付きの補足には \verb|\sidenote|\sidenote{このように表示される．番号は
章ごとに振り直される．}を，番号なしの補足には \verb|\margin| を
使う．\margin{番号なしの余白ノート．}

補足の\emph{種類}が一目でわかるように，ラベル付きの変種を3つ用意している．
\source{Amdahl の原論文~\cite{amdahl1967}，H\&P §1.9．}
\tip{試験に関係する助言や覚え方には \texttt{\string\tip} を使う．}
\caution{よくある誤解には \texttt{\string\caution} を使う．}
\begin{itemize}
  \item \verb|\source{...}| --- 出典（講義，教科書の節，論文）
  \item \verb|\tip{...}| --- 実用的な助言
  \item \verb|\caution{...}| --- 落とし穴・誤解
\end{itemize}

\section{用語と索引}

\term[ぱいぷらいん]{パイプライン}[pipeline]のような用語は \verb|\term| で
導入する．語を強調し，索引に登録し，必要なら余白に対訳（ここでは英語）を
表示する．日本語の索引は読みで並べ替えるため，漢字を含む語には
\verb|\term[よみ]{語}| のように読みを与える（カタカナ語は省略できる）．
2回目以降の\term*{パイプライン}には \verb|\term*| を使い，索引にページ番号が
並びすぎないようにする．\term[こうそくか]{高速化率}[speedup]と
\term[あむだーるのほうそく]{Amdahl の法則}[Amdahl's law]も索引に載る．

\begin{definition}[高速化率]
  実行時間のうち割合 $f$ の部分を $s$ 倍速くしたとき，全体の高速化率は
  \begin{equation}
    S = \frac{1}{(1-f) + f/s}
    \label{eq:amdahl}
  \end{equation}
  で与えられる．
\end{definition}

\begin{example}
  $f = 0.9$，$s \to \infty$ のとき，\cref{eq:amdahl} より $S = 10$ となる．
  手を付けなかった \SI{10}{\percent} が達成可能な高速化の上限を決める．
\end{example}

\begin{note}
  \verb|definition|，\verb|example|，\verb|note| は章ごとに1つのカウンタを
  共有するので，\verb|\cref| で相互参照できる．
\end{note}

\begin{keypoint}[Amdahl の法則]
  頻出ケースを最適化せよ．手を付けなかった部分が利得の上限を決める．
  余白マクロはこのボックスの中でも使える．\margin{要点ボックスの中から．}
\end{keypoint}

\section{図・表・コード}

\begin{marginfigure}
  \begin{tikzpicture}[scale=0.6]
    \draw[->] (0,0) -- (4,0) node[right,font=\tiny] {$f$};
    \draw[->] (0,0) -- (0,3) node[above,font=\tiny] {$S$};
    \draw[thick,ln-accent] plot[domain=0:3.6,samples=40] (\x,{1/(1-\x/4+0.001)*0.7});
  \end{tikzpicture}
  \caption{余白に置いた図（\texttt{marginfigure}）．}
  \label{fig:margin}
\end{marginfigure}

小さな図は余白に置く（\cref{fig:margin}）．通常の図は標準の \verb|figure|
環境を，幅の広い図は本文と余白にまたがる \verb|widefigure| を使う
（\cref{fig:wide}）．日英で共有する図は \texttt{figures/} に置き，
ラベルは \verb|\iflangja| で切り替える．

\begin{figure}[htbp]
  \centering
  % Shared TikZ figure.  Labels switch with the document language.
% Usage: % Shared TikZ figure.  Labels switch with the document language.
% Usage: % Shared TikZ figure.  Labels switch with the document language.
% Usage: \input{figures/pipeline-stages}
\begin{tikzpicture}[
  stage/.style={draw, rounded corners=1pt, minimum width=1.6cm, minimum height=0.8cm,
                font=\sffamily\small, fill=ln-box},
  >={Stealth[length=2mm]}, node distance=4mm,
]
  \node[stage] (if)  {IF};
  \node[stage, right=of if]  (id)  {ID};
  \node[stage, right=of id]  (ex)  {EX};
  \node[stage, right=of ex]  (mem) {MEM};
  \node[stage, right=of mem] (wb)  {WB};
  \draw[->] (if) -- (id);
  \draw[->] (id) -- (ex);
  \draw[->] (ex) -- (mem);
  \draw[->] (mem) -- (wb);
  \node[below=1mm of if,  font=\footnotesize\color{ln-muted}] {\iflangja{命令フェッチ}{fetch}};
  \node[below=1mm of id,  font=\footnotesize\color{ln-muted}] {\iflangja{デコード}{decode}};
  \node[below=1mm of ex,  font=\footnotesize\color{ln-muted}] {\iflangja{実行}{execute}};
  \node[below=1mm of mem, font=\footnotesize\color{ln-muted}] {\iflangja{メモリ}{memory}};
  \node[below=1mm of wb,  font=\footnotesize\color{ln-muted}] {\iflangja{書き戻し}{write back}};
\end{tikzpicture}

\begin{tikzpicture}[
  stage/.style={draw, rounded corners=1pt, minimum width=1.6cm, minimum height=0.8cm,
                font=\sffamily\small, fill=ln-box},
  >={Stealth[length=2mm]}, node distance=4mm,
]
  \node[stage] (if)  {IF};
  \node[stage, right=of if]  (id)  {ID};
  \node[stage, right=of id]  (ex)  {EX};
  \node[stage, right=of ex]  (mem) {MEM};
  \node[stage, right=of mem] (wb)  {WB};
  \draw[->] (if) -- (id);
  \draw[->] (id) -- (ex);
  \draw[->] (ex) -- (mem);
  \draw[->] (mem) -- (wb);
  \node[below=1mm of if,  font=\footnotesize\color{ln-muted}] {\iflangja{命令フェッチ}{fetch}};
  \node[below=1mm of id,  font=\footnotesize\color{ln-muted}] {\iflangja{デコード}{decode}};
  \node[below=1mm of ex,  font=\footnotesize\color{ln-muted}] {\iflangja{実行}{execute}};
  \node[below=1mm of mem, font=\footnotesize\color{ln-muted}] {\iflangja{メモリ}{memory}};
  \node[below=1mm of wb,  font=\footnotesize\color{ln-muted}] {\iflangja{書き戻し}{write back}};
\end{tikzpicture}

\begin{tikzpicture}[
  stage/.style={draw, rounded corners=1pt, minimum width=1.6cm, minimum height=0.8cm,
                font=\sffamily\small, fill=ln-box},
  >={Stealth[length=2mm]}, node distance=4mm,
]
  \node[stage] (if)  {IF};
  \node[stage, right=of if]  (id)  {ID};
  \node[stage, right=of id]  (ex)  {EX};
  \node[stage, right=of ex]  (mem) {MEM};
  \node[stage, right=of mem] (wb)  {WB};
  \draw[->] (if) -- (id);
  \draw[->] (id) -- (ex);
  \draw[->] (ex) -- (mem);
  \draw[->] (mem) -- (wb);
  \node[below=1mm of if,  font=\footnotesize\color{ln-muted}] {\iflangja{命令フェッチ}{fetch}};
  \node[below=1mm of id,  font=\footnotesize\color{ln-muted}] {\iflangja{デコード}{decode}};
  \node[below=1mm of ex,  font=\footnotesize\color{ln-muted}] {\iflangja{実行}{execute}};
  \node[below=1mm of mem, font=\footnotesize\color{ln-muted}] {\iflangja{メモリ}{memory}};
  \node[below=1mm of wb,  font=\footnotesize\color{ln-muted}] {\iflangja{書き戻し}{write back}};
\end{tikzpicture}

  \caption{古典的な5段パイプライン（共有 TikZ ソース）．}
  \label{fig:pipeline}
\end{figure}

\begin{widefigure}[htbp]
  \begin{tabular}{@{}lccccccc@{}}
    \toprule
    サイクル & 1 & 2 & 3 & 4 & 5 & 6 & 7 \\
    \midrule
    \texttt{add}  & IF & ID & EX & MEM & WB &    &    \\
    \texttt{sub}  &    & IF & ID & EX  & MEM & WB &   \\
    \texttt{lw}   &    &    & IF & ID  & EX  & MEM & WB \\
    \bottomrule
  \end{tabular}
  \caption{本文と余白にまたがる幅広フロート（\texttt{widefigure}）．}
  \label{fig:wide}
\end{widefigure}

コードには \verb|listings| を使う．
\begin{lstlisting}[language={[x86masm]Assembler}, caption={アセンブリのループ．}]
loop:   add   r1, r1, r2      ; accumulate
        addi  r3, r3, -1
        bne   r3, r0, loop    ; branch if not done
\end{lstlisting}

擬似コードには \verb|algorithmicx| を使う．
\begin{algorithm}[htbp]
  \caption{2ビット飽和カウンタの更新}
  \begin{algorithmic}[1]
    \If{taken}
      \State $c \gets \min(c+1, 3)$
    \Else
      \State $c \gets \max(c-1, 0)$
    \EndIf
  \end{algorithmic}
\end{algorithm}

\section{二言語での執筆}

両言語版は \texttt{format/}，\texttt{figures/}，\texttt{references.bib} を
共有し，\texttt{en/} と \texttt{ja/} 以下の本文だけが異なる．共有ファイルの
中では \verb|\iflangja{日本語}{英語}| か，省略形の \verb|\ja{...}|，
\verb|\en{...}| を使う．

\end{document}
